\begin{proposition}
Recall that $\FHG$ is equipped with two filtrations \textup{(}the moving and transposition filtrations\textup{)}. If we take the leading order term of $\Psi(m_{\bs\mu})$ with respect to the transposition filtration, and then the leading order term with respect to the  moving filtration, we get $X_{\bm\hat{\bs\mu}}$.
\end{proposition}

\begin{proof}
Consider $ev_n(m_{\bs\mu})$. In the notation of Lemmas \ref{cycle_label_compatibility} and \ref{cycle_label_product},
\[
L_j(c) = \sum_{i<j} X_{(i,j)}(c).
\]
Analogously to the proof of Proposition \ref{murphy_prop}, suppose that $\sigma$ is a cycle containing the index $j$. Multiplying $X_{\sigma}(1)$ by $X_{(i,j)}(c)$ inserts $i$ into the cycle if $i$ is not already in $\sigma$, or removes if it is. The size with respect to the transposition filtration is increased by 1 in the first case, and reduced by 1 in the second case. To find the leading order term of $L_j(1)^r$, we view it as a sum of products of $X_{(i_k,j)}(1)$ ($k=1,\ldots,r$) where each $i_k$ is less than $j$. The maximal degree with respect to the transposition filtration is achieved when the $i_k$ are distinct. In that case, Lemmas \ref{cycle_label_compatibility} and \ref{cycle_label_product} imply that the product of the $X_{(i_k,j)}(1)$ is $X_{\sigma}(1)$, where $\sigma = (i_r, i_{r-1}, \ldots, i_1, j)$ is an $(r+1)$-cycle, and furthermore the leading term of $L_j(1)^r c^{(j)}$ in the transposition filtration is the sum of $X_{\sigma}(c)$ across all $(r+1)$-cycles $\sigma$ whose largest element is $j$. Now, $ev_n(m_{\bs\mu})$ is a sum of products of $L_j(1)^r c^{(j)}$ where the pair of parameters $(r, c)$ occurs $m_r(\bs\mu(c))$ times. Each such term gives us an $(r+1)$-cycle of type $c$.

We maximise the degree in the moving filtration when all these cycles are disjoint. In that case, the result is an element of fully-reduced cycle type $\bs\mu$. Every such element arises exactly once, because we may reconstruct the monomial which gave rise to it: an $(r+1)$ cycle whose largest element is $j$ and has type $c$ must have come from $ev_n(x_j^r(c))$.
\end{proof}
